% ---- abstract.tex ----
% Data sources: /home/sharaths/projects/Prabh\=asa-integration/research/memory/findings.md (F1/F2/F3)
% /home/sharaths/projects/Prabh\=asa-integration/site/src/data/results.json (baseline, submission results)
% /home/sharaths/projects/Prabh\=asa-integration/data/hf_export/prabhasa_b_ss_0.1_glue/glue_summary.json (GLUE 0.5807)

Data-efficient pretraining of small language models remains an open problem, with limited evidence that linguistically grounded inductive biases improve downstream performance under fixed training budgets. This paper evaluates Prabhāsa, a controlled experimental framework examining whether Pāṇinian morphosyntactic structure, specifically through morpheme-boundary embeddings and role-aware masking strategies, yields measurable gains on the BabyLM official suite across 100M and 10M parameter models. Experiments employ a pre-registered hypothesis hierarchy, testing both masking-distribution mechanisms and auxiliary objective supervision with matched baselines, three or more seeds per arm, and paired bootstrap inference to control family-wise error. Three findings emerge. First, architecture (RoPE positional encoding) and objective choice (pure MLM without CLM) drive scale-dependent improvements: the 100M model achieves 73.06 BLiMP and 55.99 text-average, compared to the Strict-track baseline of 74.53 BLiMP, while the 10M hybrid model reaches 64.09 BLiMP compared to the Strict-Small baseline of 65.08 BLiMP. Second, kāraka role-stratified masking shows no significant causal effect at matched budget (F2 produces delta +0.10 95\% CI [−0.99, +1.20]). Third, a supervised kāraka-role-prediction auxiliary objective shows no statistically significant effect across five seeds (delta +0.76 95\% CI [−0.74, +2.27]). The contribution of the Pāṇinian framing lies in transparent mechanism design and interpretability rather than empirical lift on this benchmark. These results are reproducible, openly released on HuggingFace and GitHub, and provide a rigorous baseline for future work on linguistically motivated pretraining.
